% ============================================================
%  QSP Clock Protocol  —  Supplemental Material
%  PRA-style
% ============================================================
\documentclass[aps,pra,twocolumn,superscriptaddress,longbibliography]{revtex4-2}

\usepackage{amsmath,amssymb,mathtools}
\usepackage{graphicx}
\usepackage{booktabs}
\usepackage{bm}
\usepackage{physics}
\usepackage{xcolor}
\usepackage{xr}
\usepackage{hyperref}
\usepackage{cleveref}
\usepackage{float}
\usepackage{algorithm}
\usepackage{algpseudocode}

\graphicspath{{figures/qubit_performance_plots/}}

% Cross-reference main paper equations (compile main_pra first).
% Prefix M- avoids label collisions (e.g. eq:phop exists in both documents).
\externaldocument[M-]{main_pra}

\newcommand{\bb}[1]{\mathbf{#1}}
\newcommand{\id}{\mathbb{I}}
\newcommand{\w}{\omega}
\newcommand{\rabi}{\Omega}
\newcommand{\be}{\begin{equation}}
\newcommand{\ee}{\end{equation}}

\begin{document}

\title{Supplemental Material:\\
Noise-Resilient Atomic Clocks via Quantum Signal Processing:\\
Sensitivity, Filtering, and Multilevel Spectroscopy}

\author{Jasmine Sinanan-Singh}
\affiliation{Department of Physics, Institution A}
\author{Leon Zaporski}
\affiliation{Department of Physics, Institution B}
\author{Vladan Vuletic}
\affiliation{Department of Physics, Institution B}
\author{Isaac L. Chuang}
\affiliation{Department of Physics, Institution B}

\date{\today}

\maketitle

This Supplemental Material provides extended derivations, additional analytical results,
numerical simulation details, and pulse optimization descriptions that support the main text.
Section labels prefixed with ``S'' correspond to supplemental sections.

\setcounter{section}{0}
\renewcommand{\thesection}{S\Roman{section}}
\renewcommand{\theequation}{S\arabic{equation}}
\setcounter{equation}{0}
\renewcommand{\thefigure}{S\arabic{figure}}
\setcounter{figure}{0}


% ============================================================
\section{GPS Error Signal and DC Sensitivity for the Three-Level Clock}
\label{supp:gps_sensitivity}
% ============================================================

This section derives the exact GPS clock signal for the three-level Lambda system
and establishes the key result that all GPS-$m$ protocols share the same signal slope as Ramsey.

\subsection*{SI.1 Three-level propagator}

The Lambda system has levels $|g\rangle$, $|e\rangle$, $|m\rangle$ with probe
drive $H_{\rm drive} = \tfrac{\Omega}{2}(|g\rangle\langle e|+|e\rangle\langle g|)$
and detuning Hamiltonian $H_\delta = \delta|e\rangle\langle e|$.
The physically correct form $H_\delta = \delta|e\rangle\langle e|$ is \emph{not}
traceless in the $\{g,e\}$ probe subspace; it generates a global probe phase
$e^{-i\delta\tau/2}$ relative to $|m\rangle$.

Over a single continuous drive of duration $T$ at detuning $\delta$, the three-level
propagator is
\be
  U(\delta, T) = e^{-i\delta T/2}
  \begin{pmatrix}
    f(\delta,T) & ig(\delta,T) & 0 \\
    ig^*(\delta,T) & f^*(\delta,T) & 0 \\
    0 & 0 & e^{i\delta T/2}
  \end{pmatrix},
  \label{eq:3level_prop}
\ee
where $f$, $g$ are the amplitudes for the SU(2) rotation in the probe
subspace under the traceless Hamiltonian
$\tilde{H} = \tfrac{\Omega}{2}\sigma_x^{ge} - \tfrac{\delta}{2}\sigma_z^{ge}$:
\begin{align}
  f(\delta,T) &= \cos\!\left(\frac{\Omega_R T}{2}\right)
    + i\frac{\delta}{\Omega_R}\sin\!\left(\frac{\Omega_R T}{2}\right), \\
  g(\delta,T) &= -i\frac{\Omega}{\Omega_R}\sin\!\left(\frac{\Omega_R T}{2}\right),
\end{align}
with $\Omega_R = \sqrt{\Omega^2 + \delta^2}$.

\subsection*{SI.2 GPS clock signal}

The initial clock state is $|\psi_0\rangle = (|g\rangle+|m\rangle)/\sqrt{2}$ and
the measurement operator is $M = \sigma_y^{gm} = -i|g\rangle\langle m|+i|m\rangle\langle g|$.
After the GPS drive, $|g\rangle$ evolves to $U|g\rangle$ while $|m\rangle$ is
unaffected.  The clock signal is:
\begin{align}
  S(\delta)
  &= \langle\psi_0|U^\dagger M\,U|\psi_0\rangle
   = 2\,\mathrm{Im}\!\left[\left(U^\dagger|g\rangle\right)^* \cdot
     \frac{1}{\sqrt{2}}\right] \cdot \frac{1}{\sqrt{2}} \\
  &= \mathrm{Im}\!\left[e^{i\delta T/2}f^*(\delta,T)\right].
\end{align}
Substituting the explicit form of $f$:
\be
  \boxed{
  S(\delta) = \sin\!\left(\frac{\delta T}{2}\right)\cos\!\left(\frac{\Omega_R T}{2}\right)
            - \frac{\delta}{\Omega_R}\cos\!\left(\frac{\delta T}{2}\right)
              \sin\!\left(\frac{\Omega_R T}{2}\right).}
  \label{eq:gps_signal_supp}
\ee
This result correctly accounts for the global probe phase; the alternative
probe-subspace-traceless formula $S = (\delta/\Omega_R)\sin(\Omega_R T/2)$ omits
this phase and gives incorrect predictions away from $\delta = 0$.

\subsection*{SI.3 Signal Slope: all GPS-$m$ match Ramsey}

Differentiating Eq.~\eqref{eq:gps_signal_supp} with respect to $\delta$ and
evaluating at $\delta=0$ (where $\Omega_R\to\Omega$, $\partial_\delta\Omega_R\to 0$):
\begin{align}
  \partial_\delta S\big|_{\delta=0}
  &= \frac{T}{2}\cos\!\left(\frac{\Omega T}{2}\right)
   - \frac{\sin(\Omega T/2)}{\Omega}.
   \label{eq:gps_deriv_supp}
\end{align}
For the GPS-$m$ condition $\Omega T = 2\pi m$ (integer complete Rabi cycles):
$\sin(m\pi) = 0$ and $\cos(m\pi) = (-1)^m$, so
\be
  \partial_\delta S\big|_{\delta=0}^{\rm GPS\text{-}m} = \frac{(-1)^m T}{2},
  \qquad
  \left|\partial_\delta S\big|_{\delta=0}\right|^2 = \frac{T^2}{4}.
  \label{eq:gps_sens_supp}
\ee
The \emph{squared} signal sensitivity $T^2/4 = \pi^2$ (for $T = 2\pi$) is the same for
all $m$ and equals the Ramsey sensitivity, Eq.~(XX) of the main text.
Consequently, differences in clock performance between GPS protocols arise entirely
through the filter function $F_e(\omega)$, not through signal sensitivity.

The two-level Ramsey has double the senstivity however. If we use the Kubo variance to calculate the signal slope from the DC sensitivity of the filter function, then the signal slope $\mathcal{S} = \partial_\delta\langle M\rangle|_{\delta=0}$ can also be computed within the temporal QSP framework.
The detuning $\delta$ enters through the free-evolution phase accumulated during the pulse sequence.
Taking the derivative of $\langle M\rangle$ with respect to $\delta$ at $\delta = 0$
yields an expression structurally identical to the Kubo formula but with the deterministic
detuning Hamiltonian $\delta\sigma_z/2$ replacing the stochastic noise two-level $\beta(t)\sigma_z/2$ or three-level $\beta_e(t) \ket{e}\bra{e}$ and so to first order:

\be
  \mathcal{S} = \frac{i}{\hbar}\int_0^T dt\,
     \langle \tilde\Psi | \bigl[M_I(T),\, U_{\rm QSP}^\dagger(t)\,H_{\epsilon}\,U_{\rm QSP}(t)\bigr] | \tilde\Psi \rangle\,.
  \label{eq:slope_tqsp}
\ee
\subsection*{SI.4 GPS filter function}

Because $|\partial_\delta S|^2$ is identical for all GPS-$m$, the FOM ranking is
determined solely by the Kubo denominator
$\mathcal{K} = \int_0^\infty F_e(\omega)S(\omega)\,d\omega/(2\pi)$.
Unlike the GPS signal, $F_e(\omega)$ varies strongly with $m$:
\begin{itemize}
  \item GPS-$m=1$ ($\Omega=1$ rad\,s$^{-1}$): $F_e$ is concentrated near $\omega=\Omega=1$,
        well below the high-pass cutoff $\omega_c = 2$ rad\,s$^{-1}$.
  \item GPS-$m=8$ ($\Omega=8$ rad\,s$^{-1}$): $F_e$ peaks near $\omega=8$,
        which is above the cutoff for high-pass noise.
\end{itemize}
Protocols with $\Omega < \omega_c$ are effectively transparent to band-limited noise,
explaining the anomalously large FOM of GPS-$m=1$ under high-pass noise.


% ============================================================
\section{Complete Kubo Variance Derivation}
\label{supp:kubo_full}
% ============================================================

This section provides the full derivation of the filter function formula used
throughout the paper, including the two-level qubit case and its three-level
Lambda clock generalization.

\subsection*{SII.1 Two-level qubit sensor}

Consider a single-qubit sensor with states $\{|g\rangle, |e\rangle\}$ and a
QSP control sequence described by the SU(2) propagator
\be
  U_{\rm QSP}(t) = \begin{pmatrix} F(t) & iG^*(t) \\ iG(t) & F^*(t) \end{pmatrix},
  \qquad |F|^2 + |G|^2 = 1.
\ee
The noise Hamiltonian is $H_\epsilon = \beta(t)\sigma_z^{ge}/2$ (symmetric detuning
fluctuation).  In the toggling frame of $U_{\rm QSP}$:
\begin{align}
  \tilde{H}(t) &= U_{\rm QSP}^\dagger H_\epsilon U_{\rm QSP} \\
  &= \frac{\beta(t)}{2}\!\begin{pmatrix}
      |F|^2-|G|^2 & -2iF^*G^* \\
      2iFG        & |G|^2-|F|^2
    \end{pmatrix}\!.
\end{align}
For initial state $|\psi_0\rangle = (|g\rangle+|e\rangle)/\sqrt{2}$ and
observable $M = \sigma_y^{ge}$, the Kubo linear-response formula gives
\begin{align}
  \langle\delta M^2\rangle
  &= \int \frac{d\omega}{2\pi}\, S_\beta(\omega)\, F_e^{\rm 2L}(\omega), \\
  F_e^{\rm 2L}(\omega) &= |\Phi(\omega)|^2,
  \qquad \Phi(\omega) = \int_0^T dt\, F^*(t)G(t)\,e^{-i\omega t}.
  \label{eq:Fe_2level}
\end{align}
The filter function $F_e^{\rm 2L} = |\Phi|^2$ is the standard result for
spin-echo and dynamical decoupling sequences~\cite{Biercuk2011}, where
$\Phi(t) = F^*(t)G(t)$ is the off-diagonal path function.
For Ramsey (two $\pi/2$ pulses of negligible duration with free evolution $T$):
$F(t)\approx 1$, $G(t)\approx 0$ except during free evolution where $G(t) = i\sin(\delta t/2)$,
giving $F_e^{\rm Ramsey} \approx 4\sin^2(\omega T/2)/\omega^2$.

\subsection*{SII.2 Three-level Lambda clock: two noise channels}

For the three-level Lambda system, we have two independent noise channels:
\begin{itemize}
  \item \emph{Probe noise}: $H_e = \beta_e(t)|e\rangle\langle e|$ with PSD $S_e(\omega)$.
  \item \emph{Clock noise}: $H_m = \beta_m(t)|m\rangle\langle m|$ with PSD $S_m(\omega)$.
\end{itemize}
The QSP drive acts only on the $\{g,e\}$ probe subspace, so $|m\rangle$ is unaffected.
The three-level propagator in the probe subspace takes the block form:
\be
  U_{\rm QSP}^{(3)} = \begin{pmatrix}
    F & iG^* & 0 \\ iG & F^* & 0 \\ 0 & 0 & 1
  \end{pmatrix}.
\ee
The toggled probe noise operator is:
\begin{align}
  \tilde{A}_e(t)
  &\equiv (U_{\rm QSP}^{(3)})^\dagger |e\rangle\langle e|\, U_{\rm QSP}^{(3)} \\
  &= |G|^2|g\rangle\langle g| + |F|^2|e\rangle\langle e|
     + iF^*G^*|g\rangle\langle e| - iFG|e\rangle\langle g|.
\end{align}
The toggled clock noise operator is $\tilde{A}_m = |m\rangle\langle m|$ (unchanged).

\subsection*{SII.3 Sensitivity trajectory and classical noise variance}

For a single shot with noise realization $\beta_e(t)$, the signal mean shifts by
\be
  \delta\langle M\rangle = -i\int_0^T dt\,\beta_e(t)\,
    \langle\psi_0|[M_I(T),\,\tilde{A}_e(t)]|\psi_0\rangle,
\ee
where $M_I(T) \equiv U^\dagger(T) M U(T)$ is the \emph{interaction-frame observable}
(the lab-frame measurement $M$ conjugated by the full propagator at time $T$),
and $\tilde{A}_e(t) = U^\dagger(t)|e\rangle\langle e|U(t)$ is the toggled probe noise.
The \emph{sensitivity trajectory} is
$r(t) \equiv \langle\psi_0|[M_I(T),\,\tilde{A}_e(t)]|\psi_0\rangle$.

\textbf{General Bloch decomposition of $M_I(T)$.}
With $M = \sigma_y^{gm}$ and the three-level propagator
(block-diagonal with SU(2) block $(F(T),G(T))$ in $\{g,e\}$ and unity on $|m\rangle$),
direct matrix multiplication gives the restriction of $M_I(T)$ to the $\{g,m\}$ subspace:
\be
  M_I(T)\big|_{gm} = m_y\,\sigma_y^{gm} + m_x\,\sigma_x^{gm},
  \label{eq:MI_bloch}
\ee
where
\be
  m_y = \mathrm{Re}[F(T)],\qquad m_x = -\mathrm{Im}[F(T)].
  \label{eq:bloch_components}
\ee
The $\sigma_z^{gm}$ component vanishes identically ($m_z = 0$) because $M = \sigma_y^{gm}$
has no diagonal part in $\{g,m\}$.
$M_I(T)$ also acquires off-diagonal matrix elements coupling $\{g,m\}$ to $|e\rangle$
whenever $G(T)\neq 0$; these contribute to the sensitivity trajectory (see below).

\textbf{Case $G(T)=0$ (Ramsey, GPS-$m$, equiangular).}
When the propagator returns to the probe subspace at $t=T$ (i.e., $G(T)=0$, so
$|F(T)|=1$), the $|e\rangle$-coupling terms in $M_I(T)$ vanish.
Evaluating $r(t)$ for $|\psi_0\rangle=(|g\rangle+|m\rangle)/\sqrt{2}$ gives:
\be
  r(t) = i\,m_y\,|G(t)|^2 = i\,m_y\,\chi(t),
  \label{eq:r_G0}
\ee
where $\chi(t)\equiv|G(t)|^2$ is the population channel path function
and $m_y = \mathrm{Re}[F(T)] = \pm 1$.
For GPS-$m$ at exact resonance ($\Omega T = 2\pi m$):
$F(T)=(-1)^m$ and hence $m_y = (-1)^m$, $m_x = 0$, so $M_I(T) = (-1)^m M$.
In particular, $M_I(T) = M$ only for even $m$ (and whenever $G(T)=0$, $F(T)=+1$).
Since $m_y^2 = 1$, the filter function $\mathcal{F}_e(\omega)=|m_y|^2|\mathcal{X}(\omega)|^2=|\mathcal{X}(\omega)|^2$
is independent of the sign of $m_y$.

\textbf{Case $G(T)\neq 0$ (general pulsed QSP).}
When $G(T)\neq 0$, $M_I(T)$ has additional matrix elements coupling $|e\rangle$
to $\{g,m\}$.  The sensitivity trajectory acquires an extra coherence-channel contribution:
\be
  r(t) = i\bigl[\mathrm{Re}[F(T)]\,\chi(t)
               + \mathrm{Re}[G(T)\,\varphi^*(t)]\bigr],
  \label{eq:r_general}
\ee
where $\varphi(t)=F^*(t)G(t)$ is the coherence-channel path function.
The filter function is then no longer simply $|\mathcal{X}|^2$ but receives
cross-contributions from $\Phi(\omega)=\mathcal{F}[\varphi]$.
For sequences satisfying $G(T)=0$ the second term vanishes and Eq.~\eqref{eq:r_G0} is recovered.
The shot-to-shot classical variance of $\langle M\rangle$ is then:
\begin{align}
  \overline{\delta\langle M\rangle^2}_{\rm noise}
  &= \int\frac{d\omega}{2\pi}\,S_e(\omega)\,\Bigl|\mathcal{F}[r(t)](\omega)\Bigr|^2 \\
  &= \int_{-\infty}^{\infty}\frac{d\omega}{2\pi}\,S_e(\omega)\,|\mathcal{X}(\omega)|^2,
\end{align}
where $\mathcal{X}(\omega) = \mathcal{F}[\chi(t)] = \int_0^T|G(t)|^2 e^{-i\omega t}\,dt$.
The DC value satisfies $|\mathcal{X}(0)|^2 = \mathcal{S}^2$ identically (the DC
consistency condition, Sec.~SII.4), so the filter function and signal slope
are not independent.

For the normalized frequency variance $\sigma^2_\nu$ (main text Eq.~\eqref{M-eq:fom}),
only the resolvable frequency band contributes:
\be
  \mathcal{V} = \int_{\omega_{\min}}^{\infty}\frac{d\omega}{2\pi}\,
    S_e(\omega)\,|\mathcal{X}(\omega)|^2\,,
  \qquad \omega_{\min} = \frac{2\pi}{T}\,,
\ee
where the lower cutoff is the Fourier resolution limit.
Noise oscillating slower than one clock cycle cannot be distinguished from a
constant offset and contributes no frequency-estimation variance.

\textit{Note.} The full quantum Kubo variance $\langle\delta M^2\rangle$ contains an
additional coherence-channel contribution $\tfrac{1}{2}|\Phi(\omega)|^2$ from
the $-\sigma_z/2$ (differential g--e phase) part of $|e\rangle\langle e|$.
This term drives off-axis quantum fluctuations in the probe coherence $F^*G$
but does not contribute to the variance of the classical signal mean $\langle M\rangle$
that a clock experiment measures.

\subsection*{SII.4 DC consistency and signal slope}

At $\omega=0$, the filter function equals:
\be
  \mathcal{F}_e(0) = |\mathcal{X}(0)|^2 = \Bigl|\int_0^T|G(t)|^2 dt\Bigr|^2.
\ee
We verify that this equals the squared signal slope
$\mathcal{S}^2 = |\partial_\delta\langle M\rangle|^2_{\delta=0}$.
Treating $\delta$ as a perturbation, the first-order signal slope is
(cf.\ Eq.~\eqref{eq:r_G0} and the Dyson expansion):
\be
  \mathcal{S} = \partial_\delta\langle M\rangle\big|_{\delta=0}
  = -i\int_0^T r(t)\,dt\,,
\ee
where $r(t) = \langle\psi_0|[M_I(T),\tilde{A}_e(t)]|\psi_0\rangle$ is the
sensitivity trajectory of Sec.~SII.3, with $M_I(T)=U^\dagger(T)MU(T)$.
For sequences with $G(T)=0$, inserting $r(t)=i\,m_y\,\chi(t)$ (Eq.~\eqref{eq:r_G0}):
$\mathcal{S} = -i\int_0^T i\,m_y\,\chi(t)\,dt = m_y\,\mathcal{X}(0)$,
and since $m_y = \mathrm{Re}[F(T)] = \pm 1$ when $|G(T)|=0$,
\be
  \mathcal{S}^2 = m_y^2\,|\mathcal{X}(0)|^2 = |\mathcal{X}(0)|^2 = \mathcal{F}_e(0)\,,
\ee
confirming
\be
  \mathcal{F}_e(0) = |\mathcal{X}(0)|^2 = \mathcal{S}^2\,.
  \label{eq:dc_consistency_supp}
\ee
This \emph{DC consistency condition} is the foundation of the normalized FOM
(main text Eq.~\eqref{M-eq:fom} and Eq.~\eqref{M-eq:dc_consistency}):
$\mathcal{F}_e(0)$ serves simultaneously as the squared signal slope (sensitivity)
\emph{and} as the filter function evaluated at zero frequency, placing DC response
and off-DC noise coupling on the same scale.
In particular, the denominator of $\sigma^2_\nu = \mathcal{V}/\mathcal{F}_e(0)$
is entirely determined by the filter function, not by an independent sensitivity calculation.
Ramsey spectroscopy saturates the Fourier bound $\mathcal{F}_e^{\rm Ramsey}(0) = T^2/4$~\cite{};
general QSP sequences may deviate from this bound.

\subsection*{SII.5 Three-level filter function (probe noise)}

Collecting the results of Secs.~SII.2--SII.4, the classical noise variance due to
probe-laser frequency noise $H_e = \beta_e(t)|e\rangle\langle e|$ is:
\be
  \overline{\langle\delta M\rangle^2}_{\rm noise}
  = \int_{-\infty}^{\infty}\frac{d\omega}{2\pi}\,
    S_{\beta_e}(\omega)\,\mathcal{F}_e(\omega)\,,
  \label{eq:full_var_supp}
\ee
with
\be
  \boxed{\mathcal{F}_e(\omega) = |\mathcal{X}(\omega)|^2\,,
  \qquad \mathcal{X}(\omega) = \int_0^T|G(t)|^2\,e^{-i\omega t}\,dt\,.}
  \label{eq:Fe_3level_supp}
\ee
This is the main result of Sec.~SII.
The filter function depends only on the population channel $\chi(t)=|G(t)|^2$
and holds for all sequences satisfying $G(T)=0$ (Ramsey, GPS-$m$, equiangular).
The $\tfrac{1}{2}|\Phi|^2$ coherence-channel contribution that appears in the
full quantum Kubo variance $\langle\delta M^2\rangle$ does not contribute to the
variance of the classical signal mean $\langle M_I(T)\rangle$ (see Sec.~SII.3).

\subsection*{SII.6 Remark: clock-transition noise}

An independent noise channel $H_m = \beta_m(t)|m\rangle\langle m|$ on the clock
transition produces an additional variance contribution.
Because $|m\rangle$ is completely decoupled from the probe drive, the
$|m\rangle$ component of the final wavefunction acquires only a stochastic phase
$\varphi_m = \int_0^T\!\beta_m(t)\,dt$, giving a protocol-independent filter function
\be
  \mathcal{F}_m(\omega) = \frac{1-\cos\omega T}{\omega^2}\,.
  \label{eq:Fm_supp}
\ee
This term is the same for all pulse sequences (it depends only on $T$, not on $F(t)$
or $G(t)$) and cancels in any protocol comparison.
We therefore restrict attention to probe noise ($\mathcal{F}_e$) throughout the main text.

\subsection*{SII.7 GPS filter function}

For a GPS sequence (continuous drive at frequency $\Omega$ for time $T$), the CK
amplitudes evolve as:
\begin{align}
  F(t) &= e^{-i\delta t/2}\!\left[\cos\!\tfrac{\Omega_R t}{2}
           + i\tfrac{\delta}{\Omega_R}\sin\!\tfrac{\Omega_R t}{2}\right], \\
  G(t) &= e^{-i\delta t/2}\!\left[-i\tfrac{\Omega}{\Omega_R}\sin\!\tfrac{\Omega_R t}{2}\right].
\end{align}
At $\delta = 0$: $F(t) = \cos(\Omega t/2)$, $G(t) = -i\sin(\Omega t/2)$, so
$\varphi(t) = F^*G = i\cos(\Omega t/2)\sin(\Omega t/2) = \tfrac{i}{2}\sin(\Omega t)$
and $|G(t)|^2 = \sin^2(\Omega t/2) = \tfrac{1}{2}[1-\cos(\Omega t)]$.
The filter function therefore has oscillatory structure at the drive frequency $\Omega$,
not at harmonics — GPS $F_e$ is peaked near $\omega = \Omega$, not nulled there.


% ============================================================
\section{Roll-Off Analysis: Smoothness and Fourier Decay}
\label{supp:rolloff}
% ============================================================

\subsection*{SIII.1 Piecewise regularity of the path functions}

The high-frequency roll-off of the filter function is controlled by the regularity
of the path functions $\varphi(t) = F^*(t)G(t)$ and $\chi(t) = |G(t)|^2$.
Two qualitatively different cases arise depending on whether pulses are
instantaneous or finite-duration.

For \emph{instantaneous} (zero-duration) pulses, $\varphi(t)$ and $\chi(t)$ are
piecewise constant between pulse times, so they have jump discontinuities at
each pulse application.
A function with a jump discontinuity has Fourier transform decaying as
$|\hat f(\omega)| \sim 1/\omega$, giving $F_e \sim \omega^{-2}$.

For \emph{finite-duration square} pulses (piecewise-constant $\Omega_j$, $\phi_j$
within each segment), the CK amplitudes $F(t)$ and $G(t)$ are continuous everywhere
because the propagator evolves continuously.
Consequently $\varphi(t)$ and $\chi(t)$ are also continuous at segment boundaries.
However, their first derivatives have jump discontinuities there, because the
Rabi frequency $\Omega$ changes abruptly.
By integration by parts once, the boundary terms vanish (continuity), leaving:
$|\hat\chi(\omega)| \leq C/\omega^2$, and hence $F_e = |\hat\chi|^2 \sim \omega^{-4}$.

\subsection*{SIII.2 Smooth envelope sequences}

Two filter-function channels contribute to $\mathcal{F}_e$:
the \emph{probe channel} $|\Phi(\omega)|^2$ with $\Phi=\mathcal{F}[F^*G]$, and
the \emph{clock channel} $|\mathcal{X}(\omega)|^2$ with $\mathcal{X}=\mathcal{F}[|G|^2]$.
For sequences satisfying $G(T)=0$ (GPS-$m$, Ramsey, equiangular) only the clock channel
contributes (Sec.~SII.2); for general QSP sequences with $G(T)\neq 0$ both channels appear.
The roll-off of each channel is controlled by the smoothness of the corresponding
path function at segment boundaries.

\textbf{Probe channel} $|\Phi|^2$:
For a pulse with smooth envelope $\Omega(t)$ satisfying $\Omega(0)=\Omega(\tau)=0$
and $k$ continuous derivatives at pulse boundaries, $F^*(t)G(t)$ inherits $k$
continuous derivatives globally.
By the Riemann--Lebesgue lemma and integration by parts $k+1$ times:
\be
  |\hat\Phi(\omega)| \leq \frac{C_k}{|\omega|^{k+1}}\,,
\ee
so $|\hat\Phi|^2 \leq C_k^2/\omega^{2(k+1)}$.

\textbf{Clock channel} $|\mathcal{X}|^2$ for GPS-$m$ (single segment, $\Theta=2\pi m$):
The integer-rotation condition pins $G(0)=G(T)=0$, so $|G(t)|^2$ vanishes at both
boundaries automatically (and its first derivative vanishes since
$\partial_t|G|^2 = \sin(\theta)\Omega/2\big|_{0,T}=0$).
For a square envelope, $\partial_t^2|G|^2\big|_0 = \Omega^2/2 \neq 0$, giving $n=2$
vanishing derivatives and $\mathcal{F}_e\sim\omega^{-4}$.
For a raised-cosine envelope with $\Omega(0)=\Omega(T)=0$, $\theta(t)\sim t^3$ near
$t=0$, so $|G|^2=\sin^2(\theta/2)\sim t^6$, giving $n=6$ and $\mathcal{F}_e\sim\omega^{-12}$.

Explicit roll-off exponents are summarized in Tables~\ref{tab:rolloff_probe} and~\ref{tab:rolloff_clock}.

\begin{table}[h]
\caption{Probe-channel ($|\Phi(\omega)|^2$) roll-off for sequences with $G(T)\neq0$.
$k$ = number of continuous derivatives of $F^*(t)G(t)$ at pulse boundaries;
$|\hat\Phi|\leq C/|\omega|^{k+1}$ gives $\mathcal{F}_e^{\rm probe}\sim\omega^{-2(k+1)}$.}
\label{tab:rolloff_probe}
\centering
\begin{tabular}{lcc}
\hline
Pulse type & $k$ & $\mathcal{F}_e^{\rm probe}$ roll-off \\
\hline
Instantaneous ($t_p \to 0$) & $-1$ & $\omega^{-2}$ \\
Square (piecewise constant) & $0$  & $\omega^{-4}$ \\
Raised-cosine envelope      & $1$  & $\omega^{-6}$ \\
$C^\infty$ envelope         & $\infty$ & faster than any power \\
\hline
\end{tabular}
\end{table}

\begin{table}[h]
\caption{Clock-channel ($|\mathcal{X}(\omega)|^2$) roll-off for GPS-$m$ sequences ($G(T)=0$,
integer Rabi cycles $\Theta=2\pi m$).
$n$ = number of vanishing derivatives of $|G(t)|^2$ at $t=0,T$;
$|\hat{\mathcal{X}}|\leq C/|\omega|^n$ gives $\mathcal{F}_e^{\rm clock}\sim\omega^{-2n}$.}
\label{tab:rolloff_clock}
\centering
\begin{tabular}{lcc}
\hline
Envelope & $n$ & $\mathcal{F}_e^{\rm clock}$ roll-off \\
\hline
Square (piecewise constant) & $2$ & $\omega^{-4}$ \\
Raised-cosine ($\Omega(0)=\Omega(T)=0$) & $6$ & $\omega^{-12}$ \\
$C^\infty$ flat-top (Planck-taper) & $\infty$ & faster than any power \\
\hline
\end{tabular}
\end{table}

\subsection*{SIII.3 Constraint: fixed pulse area}

For a smooth envelope, the constraint $\int_0^{\tau}\Omega(t)\,dt = 2\pi m$ (integer
Rabi cycles for GPS-$m$) must be maintained.  The raised-cosine envelope:
\be
  \Omega(t) = \bar\Omega\left(1 - \cos\frac{2\pi t}{\tau}\right)\,,\quad
  \bar\Omega = \frac{2\pi m}{\tau}\,,
\ee
satisfies both the boundary condition $\Omega(0)=\Omega(\tau)=0$ and the area condition.
For GPS-$m$ (clock channel), this envelope achieves $\mathcal{F}_e\sim\omega^{-12}$;
for multi-segment sequences where the probe channel dominates, it achieves $\omega^{-6}$.


% ============================================================
\section{Effective Noise Model: Instantaneous Pulse Limit}
\label{supp:effective_model}
% ============================================================

\subsection*{SIV.1 Model derivation}

In the limit $\Omega_{\rm fast} \gg \beta_{\rm rms}$ (fast strong pulses), the noise Hamiltonian
$H_\epsilon(t)$ commutes with the $\sigma_z$ free-evolution operator during the
interpulse intervals but \emph{not} during pulses.
If additionally the pulse durations $\tau_p \ll 1/\omega_c$ (pulse time much less than
the noise correlation time), the noise accumulated during the pulses is negligible.
In this double limit, the noise enters only through the phase accumulated during the
free-evolution intervals:
\be
  \delta_j = \int_{t_{j-1}^{\rm free}}^{t_j^{\rm free}} \beta(t)\,dt\,,
\ee
and the noisy propagator is:
\be
  U_{\rm QSP}^{\rm noisy} = e^{i\phi_0\sigma_x/2}\prod_{j=1}^d
    e^{i(\Delta\tau_j + \delta_j)\sigma_z/2}\,e^{i\phi_j\sigma_x/2}\,.
\ee

\subsection*{SIV.2 QSP polynomial recurrences in the instantaneous model}

The canonical QSP polynomials satisfy:
\begin{align}
  f_j &= e^{i(\Delta\tau_j+\delta_j)/2}\Bigl(\cos\tfrac{\phi_j}{2}\,f_{j-1}
         - e^{-i(\Delta\tau_j+\delta_j)}\sin\tfrac{\phi_j}{2}\,g_{j-1}^*\Bigr)\,, \\
  g_j &= e^{i(\Delta\tau_j+\delta_j)/2}\Bigl(\cos\tfrac{\phi_j}{2}\,g_{j-1}
         + e^{-i(\Delta\tau_j+\delta_j)}\sin\tfrac{\phi_j}{2}\,f_{j-1}^*\Bigr)\,,
\end{align}
with $f_0 = 1$, $g_0 = 0$.
In the absence of noise ($\delta_j \equiv 0$), these reduce to standard QSP.
Expanding to first order in $\delta_j$ recovers the Kubo result of the main text.

\subsection*{SIV.3 Dynamic range and fringe-hopping}

The instantaneous model provides a natural framework for computing the \emph{capture range}
of a clock servo.
For a QSP error signal $\langle M\rangle(\delta)$,
the error signal is approximately linear in $\delta$ for $|\delta| \lesssim \delta_{\rm lock}$
where $\delta_{\rm lock}$ is the half-width of the central fringe.

For Ramsey, the signal is approximately $-\sin(\delta T_{\rm free}/2)$, so
$\delta_{\rm lock}^{\rm Ramsey} \approx \pi/T_{\rm free} \approx \pi/T$.
For GPS-$m$ at $\delta = 0$ with the \emph{correct} Hamiltonian
$H_\delta = \delta|e\rangle\langle e|$, the exact clock signal is
(main text Eq.~\eqref{M-eq:gps_signal})
\be
  S(\delta) = \sin\!\left(\frac{\delta T}{2}\right)\cos\!\left(\frac{\Omega_R T}{2}\right)
  - \frac{\delta}{\Omega_R}\cos\!\left(\frac{\delta T}{2}\right)\sin\!\left(\frac{\Omega_R T}{2}\right).
\ee
For small $\delta$ with GPS-$m$ ($\Omega T = 2\pi m$, integer $m$):
$\cos(\Omega_R T/2) \to (-1)^m$ and $\sin(\Omega_R T/2) \to 0$ to leading order, so
$S(\delta) \approx (-1)^m \sin(\delta T/2)$.
The fringe half-width is therefore
\be
  \delta_{\rm lock}^{\rm GPS} \approx \frac{\pi}{T}\,,
\ee
the \emph{same} as Ramsey, independent of $m$.
(The claim that GPS-$m$ has a fringe $m$ times narrower arises from the
probe-subspace-traceless signal formula, which is incorrect for
$H_\delta = \delta|e\rangle\langle e|$; see main text Sec.~SI.)

The probability of fringe-hopping per cycle can be estimated from the noise statistics:
\be
  P_{\rm hop} \approx \mathrm{erfc}\!\left(\frac{\delta_{\rm lock}}{\sqrt{2}\,\sigma_\delta}\right)\,,
  \label{eq:phop}
\ee
where $\sigma_\delta = [\int_0^T S_\beta(\omega)\,d\omega/(2\pi)]^{1/2}$ is the rms noise
within the interrogation bandwidth.
Since $\delta_{\rm lock}^{\rm GPS} \approx \delta_{\rm lock}^{\rm Ramsey}$,
GPS protocols at $\delta=0$ offer the same fringe-hopping robustness as Ramsey.


% ============================================================
\section{Numerical Simulation Details}
\label{supp:numerics}
% ============================================================

\subsection*{SV.1 Filter function computation via FFT}

All filter function integrals are evaluated using the following procedure.
Given a pulse sequence of total duration $T$:
\begin{enumerate}
  \item Propagate $U_{\rm QSP}(t)$ at $N_{\rm samp}$ uniformly spaced time points
        $t_k = k\,T/N_{\rm samp}$, recording $F(t_k)$ and $G(t_k)$.
  \item Form the path function samples $\varphi_k = F^*(t_k)G(t_k)$ and
        $\chi_k = |G(t_k)|^2$.
  \item Zero-pad by a factor $p$ (pad factor) to reach $pN_{\rm samp}$ total points,
        then apply the discrete Fourier transform:
        $\Phi(\omega_j) = \sum_k \varphi_k\,e^{-2\pi ijk/(pN_{\rm samp})}$.
  \item Compute $F(\omega_j) = |\chi(\omega_j)|^2$
        at positive frequencies $\omega_j = 2\pi j/(pT)$ for $j=0,\ldots,pN_{\rm samp}/2$.
  \item The maximum frequency resolved is $\omega_{\rm max} = \pi N_{\rm samp}/T$.
\end{enumerate}
Production runs use $N_{\rm samp} = 4096$ and pad factor $p=4$, giving
$\omega_{\rm max} = 2048$~rad\,s$^{-1}$ for $T = 2\pi$.
Optimization runs use $N_{\rm samp} = 1024$, $p=2$, trading frequency resolution
for speed; the resulting objective is sufficiently smooth for convergence.

\subsection*{SV.2 Unit tests}

The simulation code passes the following unit tests:
\begin{enumerate}
  \item Three-level Ramsey filter function matches analytic formula
        $F^{\rm Ramsey}(\omega) = \sin^2(\omega T/2)/\omega^2$ to $<1\%$
        for $\omega \in [0.01/T, 100/T]$ (factor of 4 smaller than the
        two-level formula~\eqref{eq:Fe_2level} because only $|\mathcal{X}|^2$
        contributes for the three-level case).
  \item GPS-$m$ sensitivity $|\partial_\delta S|_{\delta=0}|^2 = \pi^2$ for $m = 1, 2, 4, 8$
        (all equal to Ramsey), verified by finite-difference differentiation to $<10^{-6}$.
  \item GPS-$m$ filter function amplitude at $\omega = \Omega$ is nonzero (peaked),
        confirming the absence of exact nulls at the drive frequency.
  \item Noiseless Bloch sphere trajectory for QSP sequences matches analytic
        CK polynomial predictions to relative error $< 10^{-10}$.
  \item Kubo variance $\langle\delta M^2\rangle$ computed by direct integration
        matches the FOM denominator $\int F_e(\omega)S(\omega)d\omega/(2\pi)$
        to $<0.1\%$ for white and $1/f$ noise PSDs.
\end{enumerate}

% ============================================================
\section{Pulse Optimization Details}
\label{supp:optimization}
% ============================================================

This section describes the optimization algorithms used to generate the
equiangular and pulsed QSP sequences reported in the main text.

\subsection*{SVI.1 Optimization objective}

The primary reported figure of merit is the normalized frequency variance
(main text, Eq.~\eqref{M-eq:fom}):
\be
  \sigma^2_\nu(\boldsymbol{\lambda}) = \frac{\mathcal{V}(\boldsymbol{\lambda})}{\mathcal{F}_e(0;\boldsymbol{\lambda})}\,,
  \qquad
  \mathcal{V} = \int_{\omega_{\min}}^\infty \mathcal{F}_e(\omega;\boldsymbol{\lambda})\,S(\omega)\,
    \frac{d\omega}{2\pi}\,,
  \label{eq:fom_opt}
\ee
where $\omega_{\min} = 2\pi/T$ is the Fourier resolution limit.
By the DC consistency condition~\eqref{M-eq:dc_consistency},
$\mathcal{F}_e(0;\boldsymbol{\lambda}) = |\partial_\delta\langle M\rangle|^2_{\delta=0}$,
so the denominator is the squared signal slope.
The filter function $\mathcal{F}_e(\omega) = |\mathcal{X}(\omega)|^2$
and $\boldsymbol{\lambda}$ denotes the protocol parameters.

\paragraph{Optimization objective.}
Direct minimization of $\sigma^2_\nu = \mathcal{V}/\mathcal{F}_e(0)$ can converge
to degenerate solutions in which $\mathcal{F}_e(0)\to 0$ (loss of DC sensitivity),
making the ratio trivially small.
To avoid this, the optimizer instead minimizes the \emph{Ramsey-normalized objective}:
\be
  f(\boldsymbol{\lambda}) = \frac{1}{\mathcal{F}_e(0;\boldsymbol{\lambda})}
    + \frac{\mathcal{V}(\boldsymbol{\lambda})}{\mathcal{V}_{\rm Ramsey}}\,,
  \label{eq:score_rn_supp}
\ee
where $\mathcal{V}_{\rm Ramsey} = \int_{\omega_{\min}}^\infty \mathcal{F}_e^{\rm Ramsey}(\omega)
S(\omega)\,d\omega/(2\pi)$ is the noise variance of the Ramsey protocol evaluated
under the same PSD $S(\omega)$.
Both terms equal unity at Ramsey, providing a natural scale.
The first term penalizes low DC sensitivity; the second term penalizes noise variance
on the scale set by Ramsey, preventing degenerate collapses of $\mathcal{F}_e(0)$.

All integrals are evaluated numerically via FFT (Sec.~\ref{supp:numerics}).
The final reported metric for all protocols is $\sigma^2_\nu$, not $f$.

\subsection*{SVI.2 Equiangular sequence: parameter space and algorithm}

The equiangular $N$-pulse sequence~\cite{Cummins2003} has $N$ equal-duration pulses
at drive frequency $\Omega$ and phases $\phi_j$ ($j=0,\ldots,N-1$), with interpulse
free evolution set so that the total sequence duration equals $T$.
The free parameters are:
\begin{itemize}
  \item Drive frequency $\Omega \in (0, \Omega_{\rm max}]$, where
        $\Omega_{\rm max} = \min(1.5\,\omega_c, 40)$ for band-limited noise and
        $\Omega_{\rm max} = 40$~rad\,s$^{-1}$ for unbounded spectra.
  \item Pulse phases $\phi_j \in [0, 2\pi)$, $j = 0,\ldots,N-1$.
\end{itemize}
Total parameters: $N+1$ for equiangular-$N$.

The optimization proceeds in two stages:
\begin{enumerate}
  \item \textbf{Global search}: Differential Evolution (DE) with population size 20
        and maximum 500 generations, applied to the $(N+1)$-dimensional parameter space.
        DE uses the ``best/1/bin'' mutation strategy with crossover probability 0.7.
        Multiple random restarts (12 restarts with different seeds) are used to
        mitigate convergence to local minima.
  \item \textbf{Local refinement}: L-BFGS-B applied to the best DE solution,
        with gradient estimated by central finite differences.
        Convergence tolerance: $\Delta\mathcal{C} < 10^{-9}$ and
        $|\nabla\mathcal{C}| < 10^{-7}$.
\end{enumerate}

\paragraph{White noise optimum ($N=4$).}
For white noise $S(\omega)=1$, the optimal equiangular-4 parameters found by
minimizing $f$~\eqref{eq:score_rn_supp} are:
$\Omega^* T \approx 5.56$~rad, phases $\boldsymbol{\phi}^* \approx [2.48,\; 3.38,\; 2.09,\; 3.00]$~rad,
giving $\mathcal{F}_e(0)\approx10.68$ and $\sigma^2_{\nu,w}\approx3.46\times10^{-3}$.
This corresponds to $\Omega^* \approx 5.56/(2\pi) \approx 0.88$ complete Rabi cycles
per interrogation period.

\subsection*{SVI.3 Pulsed QSP sequence: parameter space and algorithm}

The pulsed QSP sequence consists of $n$ fast rotations $R(\theta_j, \phi_j)$
(rotation angle $\theta_j$ about an in-plane axis at angle $\phi_j$) separated
by $n-1$ equal free-evolution gaps of duration $\tau_{\rm free}$.
Each rotation is implemented with a fast drive at $\Omega_{\rm fast} = 20\pi$~rad\,s$^{-1}$,
giving pulse duration $\tau_j = \theta_j/\Omega_{\rm fast}$.
The free-evolution gap is determined by the constraint that the total sequence
duration equals $T$:
\be
  \tau_{\rm free} = \frac{T - \sum_{j=1}^n \tau_j}{n-1}
    = \frac{T - \sum_j\theta_j/\Omega_{\rm fast}}{n-1}.
\ee
The free parameters are $\boldsymbol{\lambda} = (\theta_1,\ldots,\theta_n,\phi_2,\ldots,\phi_n)$,
giving $2n-1$ total parameters (the first phase $\phi_1$ is fixed to zero by rotational
symmetry).

Parameter bounds:
\begin{itemize}
  \item $\theta_j \in [0.01\pi,\; \pi]$: lower bound avoids degenerate zero-rotation
        pulses; upper bound restricts to at most one Rabi half-cycle per pulse.
  \item $\phi_j \in [0,\; 2\pi)$: unconstrained rotation phase.
  \item Feasibility: $\tau_{\rm free} > 0$ requires
        $\sum_j\theta_j < (n-1)\Omega_{\rm fast}T$, which is always satisfied
        for $\Omega_{\rm fast} T = 40\pi^2 \gg n\pi$.
\end{itemize}

The optimization uses the same two-stage DE + L-BFGS-B approach as equiangular,
with budgets scaled by the number of free parameters:

\begin{center}
\begin{tabular}{cccc}
\toprule
$n$ & Free params & DE pop & DE maxiter \\
\midrule
4  & 7  & 10 & 150 \\
8  & 15 & 10 & 100 \\
13 & 25 & 8  & 80  \\
\bottomrule
\end{tabular}
\end{center}

Multiple restarts (4 for $n=4$, 3 for $n=8$, 2 for $n=13$) are performed; the best
result across restarts is refined by L-BFGS-B.

\paragraph{Optimized parameters.}
Table~\ref{tab:qsp_params_supp} records the optimized rotation angles and axis
phases for the white-noise-optimized QSP sequences ($T=2\pi$,
$\Omega_{\rm fast}=20\pi$~rad\,s$^{-1}$).
Each noise model is optimized independently; Table~\ref{tab:qsp_params_supp} shows
the white-noise winners, which are taken from the best result across both the
$\omega_{\rm max}=50$~rad\,s$^{-1}$ and $\omega_{\rm max}=4\Omega_{\rm fast}\approx251$~rad\,s$^{-1}$
optimization runs (each sequence re-evaluated with the larger bound before comparison).
Unlike the QSP~$n{=}2$ double-$\pi$ reference, the $n{=}4,8,13$ sequences
do not share a universal structure: the optimizer exploits all available phase
and amplitude degrees of freedom to suppress broadband noise.

\begin{table*}[t]
\caption{Optimized pulsed QSP parameters (white noise, $T=2\pi$,
$\Omega_{\rm fast}=20\pi$~rad\,s$^{-1}$, $\omega_{\min}=2\pi/T$).
Listed are the best-of-two-runs sequences (see text).
Each noise model yields a distinct optimized sequence; parameters for all
noise types are provided in the supplemental code.
$\sigma^2_\nu$ values match Table~I of the main text.}
\label{tab:qsp_params_supp}
\begin{ruledtabular}
\begin{tabular}{cllll}
$n$ & $\tau_{\rm free}$~(s) & $\boldsymbol{\theta}/\pi$ & $\boldsymbol{\phi}$~(rad)
    & $\sigma^2_\nu$ (white) \\
\hline
4 & 2.063
  & $[0.414,\;0.414,\;0.445,\;0.602]$
  & $[2.989,\;2.833,\;2.062,\;3.515]$
  & $4.82\times10^{-3}$ \\[3pt]
8 & 0.881
  & $[0.358,\;0.185,\;0.327,\;0.146,$
  & $[3.381,\;3.234,\;2.904,\;4.508,$
  & $2.36\times10^{-3}$ \\
  &
  & $\phantom{[}0.149,\;0.542,\;0.176,\;0.403]$
  & $\phantom{[}0.211,\;3.174,\;2.988,\;2.695]$
  & \\[3pt]
13 & 0.506
   & $[0.402,\;1.000,\;0.185,\;0.202,\;0.170,$
   & $[2.692,\;1.679,\;2.419,\;4.405,\;2.762,$
   & $1.99\times10^{-3}$ \\
   &
   & $\phantom{[}0.029,\;0.011,\;0.266,\;0.081,\;0.147,$
   & $\phantom{[}4.798,\;3.544,\;3.240,\;3.082,\;4.005,$
   & \\
   &
   & $\phantom{[}0.712,\;0.582,\;0.335]$
   & $\phantom{[}5.411,\;3.100,\;5.083]$
   & \\
\end{tabular}
\end{ruledtabular}
\end{table*}

\paragraph{Double-$\pi$ interpretation.}
The QSP $n{=}2$ reference uses the double-$\pi$ structure
$\theta=[\pi,\pi]$, $\phi=[0,0]$, which achieves maximum DC sensitivity.
The first $\pi$-pulse maps $|g\rangle\to|e\rangle$ in the probe subspace,
placing the system entirely in the excited state for the full free-evolution duration.
The noise Hamiltonian $H_\delta = \delta|e\rangle\langle e|$ then accumulates a
complete phase $\delta\tau_{\rm free}\approx\delta T$ on $|e\rangle$,
versus the half-phase $\delta T/2$ acquired by the equatorial superposition in
standard Ramsey spectroscopy.
The final $\pi$-pulse maps $|e\rangle$ back toward $|g\rangle$ to close the sequence.
This gives $\mathcal{F}_e(0)\approx T^2$, providing $4\times$ higher DC
sensitivity than continuous-drive three-level protocols.
The optimized $n{=}4,8,13$ sequences intentionally break this double-$\pi$ structure
to redistribute filter-function weight away from low frequencies,
trading some DC sensitivity for reduced noise variance $\mathcal{V}$.


       % ============================================================
       \section{Pulse Shaping: Extended Calculations}
       \label{supp:pulseshaping_ext}
       % ============================================================

       \subsection*{SVII.1 Modified recurrences for general pulse shapes}

       For a segment $j$ with general Rabi envelope $\Omega_j(t)$ and constant phase $\phi_j$,     
       the Hamiltonian $H_j = (\Omega_j(t)/2)(\cos\phi_j\,\sigma_x+\sin\phi_j\,\sigma_y)$
       generates a single-axis rotation whose accumulated angle is the cumulative pulse area       
       \be
         h_j(t) \equiv \int_{t_{j-1}}^{t}\Omega_j(t')\,dt'\,.
       \ee
       The SU(2) propagator from $t_{j-1}$ to $t$ is
       $e^{-ih_j(t)/2\,(\cos\phi_j\,\sigma_x+\sin\phi_j\,\sigma_y)}$,
       so the amplitudes satisfy the same recurrence as for square pulses
       with $\theta_j(t)\to h_j(t)$:
       \begin{align}
         F_j(t) &= \cos\tfrac{h_j(t)}{2}\,f_{j-1} -
       e^{-i\phi_j}\sin\tfrac{h_j(t)}{2}\,g_{j-1}^*\,, \\
         G_j(t) &= \cos\tfrac{h_j(t)}{2}\,g_{j-1} +
       e^{-i\phi_j}\sin\tfrac{h_j(t)}{2}\,f_{j-1}^*\,.
       \end{align}
       The endpoint values $f_j \equiv F_j(t_j)$, $g_j \equiv G_j(t_j)$ depend only on
       the total rotation angle $\Theta_j = h_j(t_j) = \int_{t_{j-1}}^{t_j}\Omega_j\,dt$
       and the phase $\phi_j$, not on the pulse shape.

       \subsection*{SVII.2 Segment contribution $\Phi_j(\omega)$ for general envelopes}

       Expanding $F_j^*(t)G_j(t)$ using the recurrence:
       \begin{align}
         F_j^*(t) &= \cos\tfrac{h_j}{2}\,f_{j-1}^* - e^{i\phi_j}\sin\tfrac{h_j}{2}\,g_{j-1}\,,     
        \\
         G_j(t)   &= \cos\tfrac{h_j}{2}\,g_{j-1}   +
       e^{-i\phi_j}\sin\tfrac{h_j}{2}\,f_{j-1}^*\,.
       \end{align}
       Multiplying and using $\cos^2(x/2)-\sin^2(x/2)=\cos x$, $2\cos(x/2)\sin(x/2)=\sin x$,
       and defining the shorthand
       $\mathcal{A}_j \equiv e^{-i\phi_j}(f_{j-1}^*)^2 - e^{i\phi_j}g_{j-1}^2$:
       \be
         F_j^*(t)G_j(t) = \cos h_j(t)\,\varphi_{j-1}
           + \tfrac{\mathcal{A}_j}{2}\,\sin h_j(t)\,,
       \ee
       where $\varphi_{j-1}=f_{j-1}^*g_{j-1}$.  Substituting $s=t-t_{j-1}$ and integrating:
       \begin{align}
         \Phi_j(\omega) &= \int_{t_{j-1}}^{t_j}F_j^*(t)G_j(t)\,e^{-i\omega t}\,dt \\
         &= e^{-i\omega t_{j-1}}\!\left[
           \varphi_{j-1}\,\tilde{I}_c^{(j)}(\omega)
           + \tfrac{\mathcal{A}_j}{2}\,\tilde{I}_s^{(j)}(\omega)
         \right]\!,
       \end{align}
       with
       \begin{align}
         \tilde{I}_c^{(j)}(\omega) &= \int_0^{\tau_j}\cos h_j(s)\,e^{-i\omega s}\,ds\,, \\
         \tilde{I}_s^{(j)}(\omega) &= \int_0^{\tau_j}\sin h_j(s)\,e^{-i\omega s}\,ds\,,
       \end{align}
       so that $\tilde{I}_c^{(j)}+i\,\tilde{I}_s^{(j)} = \tilde{e}_j^+(\omega) \equiv
       \int_0^{\tau_j}e^{ih_j(s)-i\omega s}\,ds$.
       The full filter integral is $\Phi(\omega)=\sum_j\Phi_j(\omega)$.

       \subsection*{SVII.3 Raised-cosine: Jacobi--Anger expansion}

       For the raised-cosine envelope with total rotation $\Theta_j$ over duration $\tau_j$,       
       \be
         \Omega_j(s) = \frac{\Theta_j}{\tau_j}\!\left[1-\cos\!\frac{2\pi s}{\tau_j}\right],        
         \label{eq:RC_def}
       \ee
       the cumulative area is
       \be
         h_j(s) = \frac{\Theta_j}{\tau_j}s - \frac{\Theta_j}{2\pi}\sin\!\frac{2\pi
       s}{\tau_j}\,.
       \ee
       Writing $e^{ih_j(s)} = e^{i(\Theta_j/\tau_j)s}\cdot e^{-i(\Theta_j/2\pi)\sin(2\pi
       s/\tau_j)}$
       and applying the Jacobi--Anger expansion
       $e^{iz\sin\psi}=\sum_{n=-\infty}^{\infty}J_n(z)e^{in\psi}$
       with $z=\Theta_j/(2\pi)$ and $\psi=2\pi s/\tau_j$:
       \be
         e^{ih_j(s)} = \sum_{n=-\infty}^{\infty}J_n\!\!\left(\frac{\Theta_j}{2\pi}\right)
           \exp\!\!\left[i\!\left(\frac{\Theta_j}{2\pi}-n\right)\!\omega_R s\right]\!,
       \ee
       where $\omega_R=2\pi/\tau_j$.  Integrating term by term:
       \begin{align}
         \tilde{e}_j^+(\omega) &=
       \tau_j\sum_{n=-\infty}^{\infty}J_n\!\!\left(\frac{\Theta_j}{2\pi}\right)
           \frac{\sin\xi_n^{(j)}}{\xi_n^{(j)}}\,e^{i\xi_n^{(j)}}\,,
         \label{eq:eplus_exact}
       \end{align}
       where $\xi_n^{(j)} = \bigl[(\Theta_j/2\pi - n)\omega_R - \omega\bigr]\tau_j/2$,
       with $\sin\xi/\xi \equiv 1$ at $\xi=0$.
       The real and imaginary parts give $\tilde{I}_c^{(j)}$ and $\tilde{I}_s^{(j)}$
       respectively.

       For a full half-cycle ($\Theta_j=\pi$): $J_n(1/2)$ with dominant terms
       $J_0(1/2)\approx 0.938$ (peak at $\omega=\omega_R/2$) and
       $J_{\pm 1}(1/2)\approx \pm 0.242$ (peaks at $\omega=\mp\omega_R/2$).

       \subsection*{SVII.4 Roll-off proof: probe and clock channels}

       \paragraph{Probe channel ($|\Phi|^2\sim\omega^{-6}$).}
       The function $\sin h_j(s)$ satisfies:
       \begin{itemize}
         \item $\sin h_j(0) = 0$ and $\sin h_j(\tau_j) = \sin\Theta_j$;
         \item $\tfrac{d}{ds}\sin h_j = \Omega_j(s)\cos h_j$:
               vanishes at $s=0$ and $s=\tau_j$ since $\Omega_j(0)=\Omega_j(\tau_j)=0$;
         \item $\tfrac{d^2}{ds^2}\sin h_j = \dot\Omega_j\cos h_j - \Omega_j^2\sin h_j$:
               also vanishes at the boundaries since $\dot\Omega_j(0)=\dot\Omega_j(\tau_j)=0$      
               for the raised-cosine.
       \end{itemize}
       Applying integration by parts three times to $\tilde{I}_s^{(j)}(\omega)$, each time the     
       boundary term vanishes, leaving
       \be
         \tilde{I}_s^{(j)}(\omega) = \frac{1}{(i\omega)^3}\int_0^{\tau_j}
           \frac{d^3}{ds^3}\!\left[\sin h_j(s)\right]e^{-i\omega s}\,ds\,.
       \ee
       The third derivative $\ddot\Omega_j\cos h_j - 3\Omega_j\dot\Omega_j\sin h_j -
       \Omega_j^3\cos h_j$
       is bounded, so $|\tilde{I}_s^{(j)}(\omega)|\leq C/\omega^3$
       and $\mathcal{F}_e^{\rm probe}=\tfrac{1}{2}|\Phi|^2\sim\omega^{-6}$.
       For comparison, a square envelope has $\Omega_j(0^+)\neq 0$: only one IBP is possible,      
       giving $|\tilde{I}_s|\leq C'/\omega$ and $\mathcal{F}_e^{\rm probe}\sim\omega^{-4}$.        

       \paragraph{Clock channel for GPS-$m$ (single segment, $\Theta=2\pi m$).}
       For GPS-$m$, the integer-rotation condition $\Theta=2\pi m$ ensures
       $G(0)=G(T)=0$, so $|G(t)|^2$ vanishes at both endpoints.
       Moreover $\partial_t|G|^2\big|_{0,T}=\sin\theta\cdot\Omega/2\big|_{0,T}=0$
       (since $\sin(0)=\sin(2\pi m)=0$), so the boundary smoothness of $|G|^2$
       is governed solely by higher derivatives of $\Omega(t)$.
       For a square envelope, $\partial_t^2|G|^2\big|_0 = \Omega^2/2\neq 0$,
       giving $n=2$ vanishing derivatives and $\mathcal{F}_e\sim\omega^{-4}$.
       For a raised-cosine envelope ($\Omega(0)=\Omega(T)=0$), the cumulative area
       satisfies $\theta(t)\sim t^3$ near $t=0$, giving $|G|^2=\sin^2(\theta/2)\sim t^6$,
       hence $n=6$ vanishing derivatives and
       \be
         \mathcal{F}_e^{\rm GPS,RC}(\omega)\sim\omega^{-12}\,.
       \ee

       \paragraph{Clock channel for multi-segment sequences ($\Theta_j\neq 2\pi k$).}
       For general multi-segment sequences where individual pulse rotations $\Theta_j$
       are not integer multiples of $2\pi$, $|G(t_j)|^2=\sin^2(\Theta_j/2)\neq 0$
       at segment endpoints, creating jump discontinuities in $|G(t)|^2$.
       This gives $|\mathcal{X}(\omega)|\sim\omega^{-1}$
       and $\mathcal{F}_e^{\rm clock}\sim\omega^{-2}$, independent of pulse shaping.
       Arranging $\Theta_j=2\pi k$ (complete Rabi cycles) removes these jumps and
       restores the GPS-$m$ roll-off behavior.

       \subsection*{SVII.5 GPS $m=1$: clock-channel filter function}

       For GPS $m=1$ ($\Theta=2\pi$, square envelope, $\delta=0$):
       \begin{align}
         F(t) &= \cos\tfrac{\Omega t}{2}\,,\quad G(t) = -i\sin\tfrac{\Omega t}{2}\,, \\
         |G(t)|^2 &= \sin^2\!\tfrac{\Omega t}{2} = \tfrac{1}{2}[1-\cos(\Omega t)]\,.
       \end{align}
       The clock-channel path function $\chi(t)=|G(t)|^2$ is a finite Fourier series
       with content only at $k=0$ and $k=\pm 1$.
       Evaluating $\mathcal{X}(\omega)=\int_0^T\chi(t)e^{-i\omega t}\,dt$:
       \begin{align}
         \mathcal{X}(\omega) &= \frac{T}{2}\,\mathrm{sinc}\!\left(\frac{\omega T}{2}\right)
           e^{-i\omega T/2} \nonumber\\
           &\quad - \frac{T}{4}\!\left[\mathrm{sinc}\!\left(\frac{(\omega-\Omega)T}{2}\right)
             e^{-i(\omega-\Omega)T/2}
             + (\Omega\to{-\Omega})\right]\!,
         \label{eq:chi_GPS_sq}
       \end{align}
       where $T=2\pi/\Omega$ for GPS $m=1$ (one complete Rabi cycle).
       The filter function $\mathcal{F}_e=|\mathcal{X}(\omega)|^2$ has spectral support
       concentrated near $\omega=0$ and $\omega=\pm\Omega$; for $\omega\gg\Omega$ all
       sinc envelopes decay as $\omega^{-1}$, giving
       $\mathcal{F}_e\sim\omega^{-4}$ (consistent with the square-envelope roll-off of Sec.~SIII).
       In particular, $|\mathcal{X}(\omega)|$ is extremely small for $\omega\gtrsim 2\Omega$
       because the spectral content of $\chi(t)$ is confined to harmonics $k=0,1$;
       this is the origin of the high-pass noise suppression described in the main text.

       For the raised-cosine GPS $m=1$ envelope (Eq.~\eqref{eq:RC_def} with $\Theta_j=2\pi$),
       $|G(t)|^2=\sin^2(h(t)/2)$ inherits the $\sim t^6$ boundary behavior, and
       $\mathcal{F}_e\sim\omega^{-12}$ as established in Sec.~SIII and SVII.4.

       \subsection*{SVII.6 Gradient backpropagation through the CK recurrence}

       For gradient-based optimization of shaped-pulse sequences, the gradient
       $\partial\mathcal{C}/\partial\lambda_k$ (where $\lambda_k$ is any envelope parameter)       
       is computed via backpropagation:
       \begin{enumerate}
         \item \emph{Forward pass}: propagate $(F_j, G_j)$ via the recurrence to obtain
               $\Phi(\omega_j)$, $\mathcal{X}(\omega_j)$ at all grid frequencies.
         \item \emph{Loss}: $\mathcal{C} = \sum_j w_j S(\omega_j)[\tfrac{1}{2}|\Phi_j|^2 +
       |\mathcal{X}_j|^2]$,
               $w_j = \Delta\omega/(2\pi)$.
         \item \emph{Backward pass}: adjoint recurrence for $\partial\mathcal{C}/\partial F_j$     
               and $\partial\mathcal{C}/\partial G_j$, then chain rule to
       $\partial\mathcal{C}/\partial\lambda_k$.
       \end{enumerate}
       In the current implementation, gradients are estimated by central finite differences        
       ($\epsilon=10^{-5}$), which is adequate for L-BFGS-B convergence.


% ============================================================
\section{Two-Level Reduction}
\label{supp:two_level_reduction}
% ============================================================

Setting $|m\rangle \equiv |e\rangle$, $\beta_m \equiv \beta_e$, and taking
$M = \sigma_z^{ge}$ (so $m_z=1$, $m_x=m_y=0$) gives the standard qubit filter function:
\be
  \mathcal{F}^{(2L)}(\omega) = |\Phi(\omega)|^2\,,
\ee
where $\Phi(\omega) = \int_0^T F^*(t)G(t)\,e^{-i\omega t}\,dt$.
For a Ramsey interval ($\pi/2$--$\tau$--$\pi/2$ in the instantaneous limit),
$F^*(t)G(t) = (i/2)\sin(\omega_0\tau)\Theta(\tau<t<T-\tau)$ reduces to the standard result
$\mathcal{F}^{(2L)}_{\rm Ramsey}(\omega) = 4\sin^2(\omega T/2)/\omega^2$.


% ============================================================
\section{Filter Functions for Standard Sequences}
\label{supp:filter_functions}
% ============================================================

\subsection*{SIX.1 Ramsey sequence}

For two instantaneous $\pi/2$ pulses:
\be
  \mathcal{F}_{\rm Ramsey}(\omega) = \frac{4\sin^2(\omega T/2)}{\omega^2}\,.
\ee

\subsection*{SIX.2 CPMG sequence ($n$ pulses, spacing $\tau$)}

Instantaneous pulses:
\be
  \mathcal{F}_{\rm CPMG}(\omega) = \frac{4}{\omega^2}
    \tan^2\!\left(\frac{\omega\tau}{2}\right)\sin^2\!\left(\frac{\omega n\tau}{2}\right)\,.
\ee

\subsection*{SIX.3 GPS $m$-cycle sequence}

GPS-$m$ is a single continuous drive completing $m$ full Rabi cycles ($\Omega T = 2\pi m$),
so $G(T)=\sin(\pi m)=0$ and the filter function is entirely in the clock channel:
\be
  \mathcal{F}_e^{\rm GPS}(\omega) = |\mathcal{X}(\omega)|^2\,,\quad
  \mathcal{X}(\omega) = \int_0^T|G(t)|^2 e^{-i\omega t}\,dt\,.
\ee
At $\delta=0$, $|G(t)|^2=\sin^2(\Omega t/2)$ (square envelope), which is a two-term
Fourier series with spectral content at $k=0$ and $k=\pm 1$ only
(Sec.~\ref{supp:pulseshaping_ext}, SVII.5).
The closed-form $\mathcal{X}(\omega)$ is given by Eq.~\eqref{eq:chi_GPS_sq}.
For large $\omega$, $\mathcal{F}_e^{\rm GPS}\sim\omega^{-4}$ (square envelope)
and $\mathcal{F}_e^{\rm GPS}\sim\omega^{-12}$ (raised-cosine envelope).

\subsection*{SIX.4 Roll-off behavior and pulse smoothness}

The high-frequency decay of $\mathcal{F}_e(\omega)$ depends on which channel
contributes.  For sequences with $G(T)\neq 0$ (general pulsed QSP), both channels
are present; for sequences with $G(T)=0$ (GPS-$m$, Ramsey, equiangular) only the
clock channel $|\mathcal{X}|^2$ contributes.

The probe and clock channel roll-off exponents are given by Tables~\ref{tab:rolloff_probe}
and~\ref{tab:rolloff_clock} in Sec.~SIII.2, and summarised for GPS envelopes in
Table~III of the main text.

For standard equiangular sequences (piecewise-constant $\Omega$, $G(T)\neq 0$ in general),
$F^*(t)G(t)$ has discontinuous first derivatives at pulse boundaries ($k=0$),
giving $\omega^{-4}$ roll-off in the probe channel.


% ============================================================
\section{Temporal QSP Polynomial Recurrences}
\label{supp:recurrences}
% ============================================================

During the $j$-th segment ($t_{j-1} < t \leq t_j$, Rabi frequency $\Omega_j$,
phase $\phi_j$, duration $\tau_j$):
\begin{align}
  F_j(t) &= \cos\tfrac{\theta_j(t)}{2}\,f_{j-1} - e^{-i\phi_j}\sin\tfrac{\theta_j(t)}{2}\,g_{j-1}^*\,, \\
  G_j(t) &= \cos\tfrac{\theta_j(t)}{2}\,g_{j-1} + e^{-i\phi_j}\sin\tfrac{\theta_j(t)}{2}\,f_{j-1}^*\,,
\end{align}
with $\theta_j(t) = \Omega_j(t-t_{j-1})$, $f_0 = 1$, $g_0 = 0$.
The global amplitudes $F(t)$, $G(t)$ are assembled piecewise from these per-segment
solutions; see the main text for the full piecewise construction.


% ============================================================
\section{SU(3) Extension}
\label{supp:su3}
% ============================================================

For completeness, we record the SU(3) generalization of the Kubo variance,
relevant for future extensions to multi-qutrit systems.

The noise Hamiltonian in the Gell-Mann basis is $H(t) = \sum_a \beta_a(t)\lambda_a$ with
$[\lambda_a,\lambda_b] = 2if_{abc}\lambda_c$.
In the toggling frame of an $\mathrm{SU}(3)$ control propagator $U(t)$:
\be
  \tilde{H}(t) = \sum_{a,b}\beta_a(t)R_{ab}(t)\lambda_b\,,\quad
  R_{ab}(t) = \tfrac{1}{2}\mathrm{Tr}(\lambda_a U^\dagger(t)\lambda_b U(t))\,.
\ee
For observable $M = \sum_c m_c\lambda_c$ and independent noise channels $S_{ab} = \delta_{ab}S_a$:
\be
  \langle\delta M^2\rangle = \frac{4}{\hbar^2}\sum_a\int\frac{d\omega}{2\pi}S_a(\omega)F_a(\omega)\,,
\ee
where
\be
  F_a(\omega) = \sum_{b,c,d}\left|\int_0^T dt\,R_{ab}(t)m_c f_{bcd}\,e^{-i\omega t}\right|^2\,.
\ee
This generalizes the $|\Phi(\omega)|^2$ structure of the two-level case to the full
adjoint representation of $\mathrm{SU}(3)$.

\bibliography{ref}
\end{document}
